\chapter{多模态模型}
\begin{quote}\small
\textit{本章摘要。}本章研究图像、文本、传感器和时间序列如何在同一任务中协同。先区分共同表示空间与真正的信息融合，再讨论对比学习、跨模态注意力、模态缺失和时间对齐，最后建立面向工业数据的评价方式。关键问题不是模态越多越好，而是每个模态是否提供了可对齐、可验证的互补信息。
\end{quote}
序列、视觉和语言模型各自提供证据后，还需要判断它们是否描述同一对象和同一阶段。本章先区分共享空间与信息融合，再处理对齐误差和模态缺失。联合诊断的类别标签与预警的未来事件标签必须分别定义；维修结论只有在形成后才能作为诊断证据，不能提前进入预警模型。

把照片和振动记录放在一起之前，先核对设备编号和采集时间。若照片来自维修后，振动来自维修前，直接配成一对就可能教错模型。共同嵌入空间用于寻找匹配记录，跨模态注意力让一种输入读取另一种输入，质量门控则在照片模糊或传感器失真时调整使用程度。这几种办法需要不同的训练数据。后文都用相机与振动联合诊断来说明：多出来的信息帮了什么忙，少一路输入时又该怎么办。

\section{不同模态的共同空间}
\begin{figure}[htbp]
\centering
\begin{tikzpicture}[node distance=0.75cm and 0.9cm]
 \node[idi input, minimum width=1.8cm] (img) {图像\\$v_i$};
 \node[idi input, below=0.8cm of img, minimum width=1.8cm] (txt) {文本\\$t_i$};
 \node[idi process, right=of img, minimum width=1.8cm] (ev) {视觉编码器\\$f_v$};
 \node[idi process, right=of txt, minimum width=1.8cm] (et) {文本编码器\\$f_t$};
 \node[idi state, right=1.0cm of $(ev)!0.5!(et)$, minimum width=2.0cm] (space) {共同空间\\余弦相似度};
 \node[idi output, right=of space, minimum width=2.0cm] (loss) {对比损失\\匹配/非匹配};
 \draw[idi arrow] (img)--(ev); \draw[idi arrow] (txt)--(et);
 \draw[idi arrow] (ev)--(space); \draw[idi arrow] (et)--(space); \draw[idi arrow] (space)--(loss);
\end{tikzpicture}
\caption{双塔对比学习先把不同模态映射到共同空间，再用配对关系约束表示距离。}
\label{fig:multimodal-dual-encoder}
\end{figure}
图\ref{fig:multimodal-dual-encoder}的两路编码先独立计算，再在共同空间比较匹配关系；这种交互位置允许分别缓存表示，却未提供区域级联合推理。
CLIP 通过图文对比学习获得可比较表示，BLIP-2 以轻量查询模块连接冻结的视觉与语言模型，Large Language and Vision Assistant (LLaVA) 通过视觉指令微调学习多模态交互\cite{radford2021clip,li2023blip2,liu2023llava}。BLIP-2 中的查询变换器称为 Querying Transformer (Q-Former)。三者分别强调对齐、表示连接和指令响应，并非同一损失的规模变化。
多模态模型首先要决定“对齐”发生在哪里。CLIP 采用图像编码器$f_v$和文本编码器$f_t$构成双塔结构，经过归一化后以温度参数$\tau$计算相似度：
\begin{equation}
 s_{ij}=\frac{f_v(v_i)^{\mathsf T}f_t(t_j)}{\tau}.
\end{equation}
批内对比目标让匹配图文相似度高于非匹配组合，因此 CLIP 能把“有裂纹”“轴承外圈剥落”等文本描述映射为分类原型。它的优点是图像和文本编码可以独立缓存，缺点是双塔交互浅，难以表达复杂问答和精确区域关系。

BLIP-2 用冻结视觉编码器、冻结语言模型和轻量 Q-Former 连接二者。Q-Former 通过一组可学习查询从视觉特征中提取与文本任务相关的信息，再投影到语言模型输入空间；这样可以显著减少端到端训练成本，但视觉细节会受查询数量和投影瓶颈限制。LLaVA 类结构则通常把视觉编码器输出经过线性投影后拼接到语言模型上下文中，适合图像问答和报告生成。工业场景应把生成式描述与可定位模型分开验收：语言模型可以组织证据，不能凭文字流畅性证明缺陷确实位于某个像素区域。

对于传感器、图像和文本，常见做法是先分别得到$z_s,z_v,z_t$，再进行加权融合或跨模态注意力。若传感器序列的采样频率远高于图像帧率，可以用时间窗池化、事件 token 或 cross-attention 将高频状态压缩到视觉时间点；若某个模态临时失效，则训练时加入模态 dropout，并让质量门控$g_m$根据缺失率、噪声估计和传感器健康度动态调整贡献。多模态系统的关键不是把向量简单拼接，而是明确每种模态提供何种互补证据。

\begin{table}[htbp]
\centering
\caption{多模态模型的交互位置与适用任务}
\label{tab:multimodal-interaction}
\scriptsize
\setlength{\tabcolsep}{3pt}
\resizebox{\linewidth}{!}{%
\begin{tabular}{p{2.7cm}p{3.7cm}p{3.8cm}p{3.8cm}}
\hline
模型路线 & 交互机制 & 优势 & 工业风险与适用任务 \\
\hline
CLIP 双塔 & 图像--文本共享嵌入空间 & 编码可缓存、零样本分类方便 & 细粒度定位弱，适合筛查和检索 \\
BLIP-2 & Q-Former 从视觉特征提取查询 & 冻结大模型即可迁移 & 查询瓶颈可能丢失小缺陷，适合报告草拟 \\
LLaVA 类 & 视觉 token 投影到语言上下文 & 问答和多轮解释自然 & 生成内容需外部证据，不能代替定位器 \\
交叉注意力融合 & 查询模态读取键值模态 & 能表达局部、时间和条件关系 & 计算较重，适合联合诊断与根因分析 \\
\hline
\end{tabular}}
\end{table}

表\ref{tab:multimodal-interaction}对照共享嵌入、查询瓶颈和跨模态注意力，选型时须把编码成本与所需细粒度关系同时考虑。
以“相机图像加主轴振动诊断”为例，先由视觉编码器输出工件区域 token，由时序编码器输出若干转速同步的振动 token，再将维修记录编码为文本证据。模型可以采用两阶段决策：第一阶段用 CLIP 或轻量分类器筛选明显正常与明显异常样本，第二阶段仅对异常候选执行跨模态注意力和根因排序。训练标签应同时包含缺陷类别、严重度、发生时间和证据模态；如果只有最终维修结论，没有记录哪一帧图像或哪段频谱支持结论，模型即使获得较高准确率也难以审计。

联合损失可以写为
\begin{equation}
 \mathcal{L}=\lambda_{\rm cls}\mathcal{L}_{\rm cls}+\lambda_{\rm align}\mathcal{L}_{\rm InfoNCE}+\lambda_{\rm miss}\mathcal{L}_{\rm miss}+\lambda_{\rm rank}\mathcal{L}_{\rm rank},
\end{equation}
其中分类项约束业务输出，对比项约束共同空间，模态缺失项要求模型在删去某一路输入时仍保持可接受风险，排序项则约束根因候选的相对顺序。权重不能只按数值大小设置，应根据漏诊代价、人工复核容量和各模态的标签可信度进行敏感性分析。
图像、文本、音频和传感器序列具有不同的统计结构。多模态模型通常为每种模态建立编码器$z_m=f_m(x_m)$，再通过共享投影或跨模态注意力交互。对配对图文$(x_i,t_i)$，对比学习可最大化匹配样本相似度：
\begin{equation}
 \mathcal{L}_{\mathrm{NCE}}=-\frac{1}{N}\sum_i\log\frac{\exp(s(z_i^x,z_i^t)/\tau)}{\sum_j\exp(s(z_i^x,z_j^t)/\tau)}.
\end{equation}
这里$N$为配对样本数，$i,j=1,\ldots,N$，$z_i^x$与$z_i^v$均表示第$i$幅图像的单位范数表示，$s(a,b)=a^{\mathsf T}b$是不含温度的余弦相似度。因此$s_{ij}=s(z_i^v,z_j^t)/\tau$是已经缩放的对比得分，不能再次除以温度，且$\tau>0$。

\paragraph{融合方式。}
早期融合将原始特征拼接，晚期融合分别预测后再组合，中间融合在表示空间通过 cross-attention 交互。编码器--解码器架构还可以让一个模态条件化另一个模态的生成。工业场景中，时间同步、传感器缺失和模态质量差异往往比网络结构更决定系统效果。

\paragraph{多模态推理。}
图像问答、视觉定位和设备状态描述都可看作条件生成。应分别评估感知正确性、跨模态对齐和语言表达，避免只用语言流畅度掩盖视觉判断错误。

\paragraph{对比学习的温度参数。}
设图像和文本编码器输出归一化表示$z_i^v,z_i^t$，相似度取余弦相似度。温度$\tau$控制 softmax 的尖锐程度：
\begin{equation}
 s_{ij}=\frac{(z_i^v)^{\mathsf T}z_j^t}{\tau}.
\end{equation}
小温度会放大正负样本差异，梯度更集中于难负样本，但也可能使训练不稳定。对称的图到文、文到图损失通常比单向损失利用更多配对信息。

\paragraph{跨模态注意力。}
以文本为查询、视觉 patch 为键和值时，文本 token 可以选择性读取图像区域：

为说明其计算含义，令$Q\in\mathbb R^{n_t\times d}$、$K\in\mathbb R^{n_v\times d}$、$V\in\mathbb R^{n_v\times d_v}$。矩阵$QK^{\mathsf T}$的第$(a,b)$项是第$a$个文本查询与第$b$个视觉键的内积；除以$\sqrt d$是为了在各维独立、方差为一时保持内积方差在可控范围。对每个文本位置$a$，定义
\[
\alpha_{ab}=\frac{\exp(q_a^{\mathsf T}k_b/\sqrt d)}{\sum_{r=1}^{n_v}\exp(q_a^{\mathsf T}k_r/\sqrt d)},\qquad \sum_b\alpha_{ab}=1.
\]
于是输出第$a$行为$h_a'=\sum_b\alpha_{ab}v_b$，它是视觉值向量的凸组合，而不是一个自动产生定位真值的坐标。若缺陷区域很小，$n_v$过小或下采样会使多个区域共享一个 token；若时间掩码错误，注意力仍会在错误窗口内给出看似合理的权重。实际部署应同时检查注意力之外的区域标注、遮挡实验和跨设备稳定性。
\begin{equation}
 H_{text}'=\softmax\left(\frac{Q_{text}K_{image}^{\mathsf T}}{\sqrt d}\right)V_{image}.
\end{equation}
此处$Q_{text}\in\R^{n_t\times d}$、$K_{image}\in\R^{n_v\times d}$、$V_{image}\in\R^{n_v\times d_v}$，$n_t,n_v$分别为文本与视觉 token 数。各行存放列向量的转置，softmax 沿视觉位置逐行计算，输出为$n_t\times d_v$矩阵。
这与简单拼接不同，因为每个文本位置可使用不同的视觉证据。多层交互可以支持图像描述、区域问答和视觉定位，但其对齐质量仍受配对数据噪声影响。

\paragraph{模态缺失与质量加权。}
真实工业系统常出现相机被挡、传感器掉线或文本记录晚到。模型应显式接收“这一模态当前是否可用”的掩码$m_m$，并考虑
\begin{equation}
 z=\sum_m\frac{m_m\exp(a_m)}{\sum_r m_r\exp(a_r)}z_m.\label{eq:multimodal-quality-fusion}
\end{equation}
其中$m,r$遍历模态，$m_m\in\{0,1\}$表示第$m$个模态是否可用，$a_m$为有限的质量得分，各$z_m$须投影到相同维度。式\eqref{eq:multimodal-quality-fusion}要求$\sum_r m_r>0$；全部缺失时不计算这一比值，而应拒绝输出或调用预先验证的回退流程。训练时进行模态 dropout，能减少模型对某一模态的脆弱依赖，但必须保留真实缺失模式。

\section{时间对齐、缺失鲁棒性与工业诊断}
\paragraph{时间对齐的多模态传感器。}
视觉帧、振动窗口和控制指令可能不在同一时间尺度。可以用时间编码、可学习延迟或 cross-attention 在邻近时间段内寻找对应信息。若把未来传感器值错误地对齐到当前图像，离线结果会明显偏高，部署时却失效。因此多模态数据集必须保存采集时间、传输延迟和有效时间区间。

设第$m$种模态在时间窗$W_i=[t_i^-,t_i^+]$内产生输入$x_i^{(m)}$，编码器输出$z_i^{(m)}
\in\mathbb R^d$。所谓配对，至少应满足设备标识一致、事件窗有足够重叠，并且标签在预测时刻已经可用。若两个模态的有效窗交集为空，将它们强行作为正对会把时间错配当成监督噪声。对于存在未知延迟的系统，可在训练中显式搜索有限延迟集合$\Delta$，例如令
\[
\operatorname{sim}_i=\max_{\delta\in\Delta}\operatorname{sim}\bigl(z_i^{(v)},z_{i,\delta}^{(s)}\bigr),
\]
但最大化也可能选择偶然相似的错误事件，因此必须在留出设备和固定延迟范围上验证，而不能把对齐搜索结果直接解释为真实因果延迟。

\paragraph{图像与工艺曲线联合诊断。}
视觉编码器提取表面纹理，序列编码器提取温度和压力变化，再用融合层预测缺陷类型与严重度。消融实验至少包括仅视觉、仅时序、简单拼接和跨模态注意力四种设置。若融合模型只在完整数据上有效，还需报告相机缺失或曲线缺失时的退化曲线。

\paragraph{对齐、融合与协同的区别。}
对齐回答“不同模态是否指向同一个对象或事件”，融合回答“如何联合使用这些信息”，协同则进一步要求联合表示在任务上优于任何单一模态。三者不能混为一谈：图像和文本可以在共享空间中相近，却未必能完成精确定位；多个传感器可以简单拼接，却没有学到跨模态因果关系。

\begin{figure}[htbp]
\centering
\begin{tikzpicture}[>=Latex, node distance=0.8cm, every node/.style={font=\small}]
 \node[draw, rounded corners, fill=blue!10, minimum width=2.4cm] (vision) {视觉编码器};
 \node[draw, rounded corners, fill=orange!12, right=2cm of vision, minimum width=2.4cm] (text) {文本编码器};
 \node[draw, rounded corners, fill=green!12, below=1.1cm of vision, minimum width=2.4cm] (sensor) {传感器编码器};
 \node[draw, rounded corners, fill=yellow!20, right=2cm of sensor, minimum width=2.4cm] (fusion) {对齐/融合层};
 \node[draw, rounded corners, fill=red!10, right=2cm of fusion, minimum width=2.4cm] (out) {诊断或生成};
 \draw[->, thick] (vision)--(fusion); \draw[->, thick] (text)--(fusion); \draw[->, thick] (sensor)--(fusion); \draw[->, thick] (fusion)--(out);
\end{tikzpicture}
\caption{工业多模态系统中，编码、对齐、融合和任务输出是不同层次。}
\label{fig:multimodal-industrial-fusion}
\end{figure}
图\ref{fig:multimodal-industrial-fusion}将各模态编码器与融合层分开，说明原始输入先形成各自表示，再由对齐与融合环节为诊断或生成提供联合证据。

\paragraph{对称对比损失。}
图到文和文到图两个方向都应参与训练。令$S_{ij}=s_{ij}$为已经除以温度的得分矩阵，且$i,j=1,\ldots,N$，则对称损失为
\begin{equation}
 \mathcal{L}=\frac12\left[-\frac1N\sum_i\log\frac{e^{S_{ii}}}{\sum_j e^{S_{ij}}}
 -\frac1N\sum_i\log\frac{e^{S_{ii}}}{\sum_j e^{S_{ji}}}\right].
\end{equation}
批次中的负样本并不都是真负样本：两个不同时间的图像可能描述同一设备状态，两个文本也可能共享相同故障原因。因此对比学习需要检查 false negative，并在必要时使用组标签、软标签或难负样本过滤。

\paragraph{模态不平衡与缺失鲁棒性。}
某一模态信噪比更高时，模型可能忽略其他模态；训练损失下降并不表示融合真正发生。可以通过模态 dropout、分支独立监督、梯度平衡和质量门控减轻这一问题。测试时应绘制从完整输入到逐模态缺失的性能曲线，并区分随机缺失、连续掉线、延迟到达和系统性遮挡。

\paragraph{时间对齐与事件粒度。}
图像帧、振动窗口和控制指令的时间粒度不同。对齐时要明确事件发生时间、采集时间、传输时间和标签确认时间。窗口重叠会增加有效样本数，却也可能把同一故障事件复制到训练和测试两侧。应以事件或设备为单位划分，并在报告中说明窗口长度、步长和允许的预测提前量。

\paragraph{多模态评价。}
多模态系统应分别报告单模态基线、简单拼接、注意力融合和缺失模态结果。分类任务可检查跨模态遮挡后的鲁棒性；定位任务需要区域级准确率；生成任务需要事实一致性和证据覆盖率。最重要的消融不是把某一层删掉，而是回答“哪一种模态信息在什么工况下真正改变了决策”。

\paragraph{本章小结。}
图像、振动和文字的采集频率不同，质量也不同，还可能缺失。联合使用时，要说明每一路怎样编码、按什么时间和设备配对、相互矛盾时听谁的，以及少一路输入时还能完成什么任务。最后在独立数据上分别测试完整输入和缺失输入，才能判断复杂的融合结构是否真的有帮助。